% === ΛΩ Reflexive Equivalence inserts ===
\section{ΛΩ Reflexive Equivalence and Perfect Self–Containment}

\begin{theorem}[ΛΩ Reflexive Equivalence]
\label{thm:LambdaOmegaEquivalence}
Let the \emph{Universal Generative Principle} (UGP/GTE) define the computable substrate 
through a sequence of canonical Λ–operations, and let the \emph{Z–hypercomplex recurrence}
represent the uncomputable self–closure through Ω–operations of the form
\[
Z = \Lambda + j\,\Omega, \qquad (2\pi\Lambda)^2 + g(\Omega)^2 = (m\pi)^2,\quad m\in\mathbb{Z}_{>0}.
\]
Then $\Lambda$ and $\Omega$ are dual projections of a single SU(2) rotation in information–space:
\[
\mathcal{U}(m) = e^{\,i\,2\pi\Lambda + j\,g(\Omega)} \in SU(2),
\]
such that the Λ–axis generates the deterministic law (computable dynamics) 
and the Ω–axis enforces the reflexive closure (self–description).
Therefore the partial differential equation emerging from $\Lambda$
already satisfies its own reflexive closure under $\Omega$,
\[
\partial_t \Phi(\Lambda) = \mathcal{L}_{\Lambda}\Phi 
\;\Longleftrightarrow\;
\mathcal{R}_{\Omega}\,\Phi = 0.
\]
\end{theorem}

\begin{remark}[Discussion: Relation to Transputation and PSC]
\label{rem:TransputationRelation}
Transputation (PT) expresses the \emph{process} by which execution and evaluation coincide;
ΛΩ Equivalence expresses the \emph{structure} of that process as an SU(2) closure.
This unifies prior PSC demonstrations (thermodynamic, informational, geometric) under one algebraic symmetry.
The discrete harmonics $m=1,2,3$ act as eigen–choices of a self–referential phase (micro Λ, meso ΛΩ, macro ΛΩΛ). 
\end{remark}
